\chapter{Гейт граничной BF/AKSZ pairing form}

Квантованная симплектическая форма
\begin{equation}
 \Omega_k=\frac{k}{2\pi}\int_\Sigma\delta B\wedge\delta A,
 \qquad k\in\mathbb Z,
\end{equation}
фиксирует скобку фазовых координат, но не является положительной метрикой:
\(\Omega_k(v,v)=0\). Положительная форма требует совместимой комплексной
структуры \(J\), \(g_J(u,v)=\Omega_k(u,Jv)\).

Уже для одной канонической пары допустимо семейство
\begin{equation}
 J_s=\begin{pmatrix}0&-s^{-1}\\s&0\end{pmatrix},
 \qquad
 g_s=k\begin{pmatrix}s&0\\0&s^{-1}\end{pmatrix},
 \qquad s>0.
\end{equation}
Уровень \(k\) квантован, но поляризация \(s\) остаётся непрерывной.
Кроме того, чистое BF/AKSZ pairing кодирует интегральные
intersection/linking-данные, а объёмы \(\pi^2\), \(2\pi\) и норма
\(\pi^{-1}\) появляются только после Hodge- или metric-ввода.

Носитель \(K=\mathbb{RP}^3\times S^1\) замкнут. BFV-граница требует
разрезания или физического дефекта, то есть новых граничных данных.

\begin{nogocriterion}[Pure BF/AKSZ pairing no-go]
Квантованный топологический pairing не выводит положительную
двухсекторную trace--Hodge-метрику. Для этого требуется отдельно вывести
поляризацию, gauge fixing, гамильтониан или дефект.
\end{nogocriterion}

Следующий кандидат использует канонические условные ожидания
\begin{equation}
 E_{S^1}=\operatorname{id}\otimes\tau_{S^1},
 \qquad
 E_{\mathbb{RP}^3}=\tau_{\mathbb{RP}^3}\otimes\operatorname{id}.
\end{equation}

Машинный аудит сохранён в
\texttt{s2t\_v3\_bf\_aksz\_pairing\_results.json}.